\chapter{BPCC：身体预测控制小脑}

\begin{chapteroverviewbox}
本章讨论 Body Runtime 中最核心的快速动力学组件：BPCC（Body Predictive Control Core，身体预测控制核心）。在本书的总体架构中，BPCC 不是第二个 Agent，也不是缩小版的认知内核，更不是一个简单的运动控制器。它承担的是介于原始身体动力学与高层认知之间的一组高速、局部、连续、可适应的闭环功能：预测身体在短时间尺度上的状态变化，在给定高层控制参考与安全约束的条件下选择局部动作，估计预测与现实之间的残差，并在不破坏整体语义接口的前提下持续校准自身内部动力模型。

从多尺度智能动力学的角度看，BPCC 的存在来自一个不可回避的事实：现实身体的有效控制周期通常远短于高层认知循环。若所有运动、设备、浏览器、车辆或机器人状态都必须先进入工作记忆 $h(t)$，再经过显式规划后才能得到动作，则 Agent 的高层认知将被大量快速、重复、低语义增量的身体细节淹没，工作记忆有限容量定理也会立即成为整个具身系统的瓶颈。因此，本章把 BPCC 定义为身体尺度上的局部预测--控制--残差--适应闭环，并提出
\[
\boxed{
BPCC
=
FastPrediction
+
FastControl
+
ResidualEstimation
+
AdaptiveLearning
}
\]
这一基本结构。

本章特别坚持三个边界。第一，BPCC 可以拥有复杂的内部潜在状态，但它不拥有独立的长期身份、生命周期记忆与主体自治，因此不能被视为第二主体。第二，BPCC 的预测永远不能取得现实权威；真正的身体状态必须重新经过现实、传感和 Body Runtime 的观测链路确认。第三，BPCC 的学习主要发生在快速局部动力学与控制映射上，它不得未经治理地修改 Agent 的长期目标、自我结构和价值结构。由此，BPCC 成为本书多尺度原则在身体层的一个具体实现：\textbf{正常快速动力学留在局部闭环中解决，只有持续、重要而无法局部消解的残差才应向更高认知尺度升级。}
\end{chapteroverviewbox}


\section{为什么 Body Runtime 需要自己的智能}

如果把一个长期 Agent 看成传统的软件调用链，一个最直接的设计方式似乎是：传感器把全部状态送入认知内核，认知内核完成理解、规划和推理，然后把动作指令直接发送给执行机构。从抽象功能图上看，这条链路非常简洁：
\[
World
\rightarrow
Perception
\rightarrow
AgentKernel
\rightarrow
Action
\rightarrow
World.
\]
然而一旦 Agent 真正进入具有连续动力学的现实身体，这一结构就会暴露出根本性问题。高层认知与身体控制所处的特征时间尺度并不相同。设高层工作记忆闭环的特征时间为 $\tau_h$，身体局部控制周期为 $\tau_B$，对于多数具身系统通常要求
\[
\tau_B \ll \tau_h.
\]
在机器人姿态稳定、机械臂接触控制、车辆轨迹跟踪、数字桌面的鼠标移动、浏览器页面快速变化乃至设备资源管理等场景中，环境可能在高层推理尚未结束之前就已经发生多次有效状态转移。因此，如果认知内核必须参与所有最小控制周期，就相当于要求一个慢尺度控制器直接闭合远高于自身有效带宽的快速反馈环。这不仅增加时延，而且会产生严重的控制不稳定性。

更重要的是，这一问题并不只是传统意义上的 ``实时性不足''。在本书前面的工作记忆有限容量理论中，我们已经把 $h(t)$ 理解为有限的显式跨尺度认知界面，而不是全系统所有状态的实时镜像。若令身体每个自由度、每个传感器变化和每个短期预测都不断进入 $h(t)$，则运行态的显式结构负荷可粗略表示为
\[
C_h(t)
=
\sum_{i=1}^{N_B}
\alpha_i C_i(t),
\]
其中 $N_B$ 是被提升到高层的身体状态项数，$C_i$ 表示该状态在当前任务中的关系负荷，$\alpha_i$ 表示其进入显式认知的权重。当大量低层状态持续满足 $\alpha_i>0$ 时，
\[
C_h(t) \rightarrow C_{\max}(h),
\]
高层工作记忆就会把宝贵的关系容量消耗在本应由局部动力学自动闭合的问题上。真正成熟的智能恰恰应呈现相反方向：大量可预测、可重复、可局部修正的身体结构逐渐从显式认知中退出，只在出现异常时重新进入高层。

因此，Body Runtime 不能只是一个无状态设备驱动层。它至少需要具有局部状态估计、短期预测、动作调整和误差学习能力。这里所谓 ``自己的智能'' 并不意味着它拥有自己的自我，而是指它具有在局部状态空间中完成有效闭环所需的最小适应性。若身体状态记为 $x_B(t)$，局部控制为 $u_B(t)$，外部扰动为 $w(t)$，则身体动力学可写成
\[
\dot{x}_B
=
F_B(x_B,u_B,w).
\]
如果每次计算 $u_B$ 都必须经过 Agent Kernel，则高层系统实际承担的是
\[
u_B(t)
=
\Pi_K
\left(
h(t),E,\Theta,\Psi,\ldots
\right),
\]
这使快速身体环的带宽受限于认知环。而 BPCC 的引入使我们能够改写为
\[
u_B(t)
=
\Pi_B
\left(
z_B(t),r_K(t),\mathcal E_K(t)
\right),
\]
其中 $z_B$ 是 BPCC 的局部内部状态，$r_K$ 是高层给出的控制参考，$\mathcal E_K$ 是高层给出的允许域或预测包络。高层不再决定每一个微小动作，而是决定更慢尺度上的方向、目标与边界。

从控制论角度，这与层级控制、模型预测控制、快速内环--慢速外环系统具有明显联系；从认知科学角度，它也与熟练动作的程序化和自动化具有形式相似性。但本书并不把 BPCC 直接等同于生物小脑。``小脑''在这里首先是工程类比：它提醒我们，一个成熟智能系统必须允许大量高频预测控制在不占用显式工作记忆的情况下局部运行。BPCC 因此是 Body Runtime 从 ``传感器--执行器适配层'' 升级为 ``局部智能动力层'' 的关键节点。

\section{BPCC 不是真正主体}

BPCC 具有预测、控制、误差估计和在线学习能力，因此很容易产生一个概念误区：既然它能够观察自身局部状态、预测未来并改变动作，它是否已经构成了一个独立 Agent？在本书的架构中，答案是否定的。我们必须区分 ``局部闭环智能'' 与 ``持续主体''。局部闭环能够表现出高度复杂和自适应行为，但主体性的工程定义还要求更慢尺度上的身份连续、生命周期记忆、目标整合、自我模型以及跨任务一致性。BPCC 并不拥有这些结构。

设一个完整 Agent 的生命周期状态近似为
\[
\mathcal X_A(t)
=
\left[
h(t),E(t),\Theta(t),\Psi(t),\Omega(t),G(t),\ldots
\right].
\]
其中 $\Psi$ 承担长期自我与身份连续结构，$\Omega$ 表示更慢的生命周期生成倾向，$E$ 与 $\Theta$ 则保存跨任务和跨时间的经验与结构记忆。与此不同，BPCC 的内部状态可以写成
\[
z_B(t)
\in
\mathcal Z_B,
\]
它主要编码当前身体动力学、短期可预测状态、局部控制历史和残差统计。其有效记忆范围通常满足
\[
\tau_{z_B}
\ll
\tau_E,\tau_\Theta,\tau_\Psi.
\]
因此，BPCC 可以 ``记住'' 最近几十个控制周期发生了什么，却不因这种局部记忆而获得 ``我是谁'' 的生命周期身份。

这一差别可以进一步从目标来源上表达。主体级目标通常来自
\[
G_A
=
G(h,E,\Theta,\Psi,\Omega,\text{World}),
\]
而 BPCC 的局部优化目标则由高层目标经过身体语义和约束映射后给出：
\[
r_B
=
\Gamma_B(G_A,\text{BodyState},\text{Context}).
\]
BPCC 可以在局部动作空间中选择最优控制，但它不应自行把 ``保持机械臂末端轨迹稳定'' 重写成 ``改变整个 Agent 的人生目标''。也就是说：
\[
\boxed{
BPCC\;controls\;execution,
\quad
but\;does\;not\;own\;identity.
}
\]

从因果权限的角度，我们可以把系统变量分成不同写权限集合。若 $\mathcal V_B$ 表示身体快速变量，$\mathcal V_K$ 表示认知慢变量，则 BPCC 的直接写入权限应近似满足
\[
Write(BPCC)
\subseteq
\mathcal V_B,
\]
而对于
\[
\Psi,\Omega,\Theta_{\mathrm{core}}
\]
等长期认知变量，应有
\[
Write_{direct}(BPCC,\Psi)
=
Write_{direct}(BPCC,\Omega)
=
0.
\]
BPCC 可以通过长期残差间接影响高层学习，但这种影响必须经由可观察的残差信号和上层治理链路发生，而不是通过隐藏状态直接修改主体结构。

这一设计与本书此前的 ``单一持续主体原则'' 一致。我们允许系统内部存在许多局部自治环：记忆检索环、工具调用环、身体控制环、环境控制环，甚至外部执行进程；但这些闭环并不会自动获得新的主体身份。否则，一个长期 Agent 会因为每个高性能子系统的存在而被人为拆分成许多互相竞争的 ``小 Agent''，这不仅破坏身份连续性，也会带来权限冲突和控制竞争。

因此，本章坚持：
\[
\boxed{
LocalAutonomy
\neq
SubjectAutonomy.
}
\]
BPCC 可以拥有高度自治的快速局部控制，但这种自治发生在由主体给出的可行域内。这一结构使我们能够在不牺牲单一主体连续性的前提下获得复杂的分层智能，也使 Body Runtime 可以随着身体越来越复杂而继续扩展，而不必把每一个新增身体闭环都升级成新的 Agent。

\section{Fast Prediction：短时间尺度的身体预测}

BPCC 的第一项核心能力是快速预测。身体控制如果完全依赖当前观测，控制器实际上只能在扰动已经发生之后被动修正；而具有模型的系统可以提前估计未来几个控制周期中的状态变化，在风险真正出现之前修正动作。BPCC 的预测不是关于人生、任务或者复杂社会未来的高层想象，而是一个高度局部化的身体动力预测问题。

令身体可观测状态为
\[
x_t^B,
\]
控制输入为
\[
u_t,
\]
外部扰动为
\[
w_t.
\]
一个基本离散动力模型可以写成
\[
x_{t+1}^B
=
f_\phi(x_t^B,u_t,w_t)
+
\epsilon_t,
\]
其中 $f_\phi$ 是 BPCC 当前学习到的局部动力模型，$\phi$ 为其可适应参数，$\epsilon_t$ 表示未被模型解释的真实扰动。BPCC 的任务不是只预测一个点，而应形成短时间预测轨迹：
\[
\hat{X}_{t:t+H}
=
\left\{
\hat x_{t+1},
\hat x_{t+2},
\ldots,
\hat x_{t+H}
\right\}.
\]
这里的预测长度 $H$ 由身体动力学的时间尺度决定，不应机械追求越长越好。高频局部模型预测得越远，模型误差累积和环境不确定性通常越明显。因此 BPCC 更关注一个有限预测窗口：
\[
H_B\Delta t
\ll
\tau_h.
\]

这一短视预测具有一个重要作用：它把高层抽象意图转换成身体尺度的可达性判断。例如，认知内核可能只给出 ``把手移动到目标区域'' 的高层要求，而 BPCC 必须判断在当前关节速度、惯量、负载、接触状态和局部障碍条件下，未来几百毫秒哪些轨迹仍然可行。于是高层意图 $g_K$ 并不直接变成电机控制，而是先进入身体可行性空间：
\[
\mathcal X_{feasible}^{B}(t:t+H)
=
\left\{
X:
C_B(X,U)\leq 0
\right\}.
\]
BPCC 的预测因此既是状态预测，也是局部可达性预测。

为了使预测具有工程价值，我们还必须为其引入不确定性。若只给出
\[
\hat x_{t+1},
\]
而没有预测置信度，则控制器可能把脆弱模型的输出误当成现实。更合理的形式是：
\[
p(x_{t+1:t+H}
\mid
x_t,u_{t:t+H-1})
\]
或者用均值--协方差近似：
\[
\hat x_{t+k}
\sim
\mathcal N(\mu_{t+k},\Sigma_{t+k}).
\]
这样 BPCC 可以根据预测不确定性自动收缩动作幅度。例如定义局部控制激进度
\[
\beta_t
=
\beta_0
\exp(-\lambda\,\mathrm{Tr}(\Sigma_t)),
\]
当模型不确定性上升时，控制策略自然趋向保守。

Fast Prediction 还承担一个更深的认知卸载作用。若一种身体过程长期可预测，那么高层 Agent 就没有必要持续观察每一个细节。我们可以把 ``身体过程是否需要上浮'' 写成：
\[
\Gamma_B(t)
=
F(
\|\varepsilon_B(t)\|,
R_B(t),
U_B(t),
P_B(t)
),
\]
其中 $\varepsilon_B$ 为预测残差，$R_B$ 为风险，$U_B$ 为不确定性，$P_B$ 为异常持续性。只有当
\[
\Gamma_B(t)\geq\theta_B
\]
时，局部过程才需要被转换为更高层语义事件。这是本书 ``Reality does not continuously become thought; relevant deviation becomes thought'' 原则在身体层的具体实现。

因此，Fast Prediction 的真正价值并不只是 ``让动作更准''。它让 Body Runtime 获得一个能够自行消化可预测现实的时间窗口，从而把高层认知从大量短周期反馈中解放出来。BPCC 预测得越稳定，Agent 就越能够把显式认知资源保留给新颖、不确定、具有长期影响的问题。

\section{Fast Control：在高层约束下闭合局部动作}

如果预测回答的是 ``接下来可能发生什么''，Fast Control 回答的则是 ``在有限时间窗口中现在应该怎样改变身体状态''。BPCC 的控制目标并不是自由生成任意行为，而是在高层给出的目标、边界和安全约束下寻找局部可执行控制。我们因此把高层与 BPCC 的关系表示为：
\[
\boxed{
HigherLevelSetsFeasibleRegion,
\qquad
LowerLevelOptimizesInsideIt.
}
\]

假设认知内核给出的高层控制参考为 $r_t$。这一参考可以是目标位姿、速度区间、浏览器中的目标 UI 元素、设备状态目标、导航方向或更抽象的身体意图。BPCC 不应把 $r_t$ 理解成需要逐字执行的动作脚本，而应把它映射成一个局部优化问题。最基本形式可写成：
\[
u_t^\ast
=
\arg\min_{u_t}
\mathcal L_B
(x_t,u_t,r_t)
\]
subject to
\[
x_{t+1}=f_\phi(x_t,u_t),
\qquad
g_B(x_t,u_t)\leq0.
\]
这里 $\mathcal L_B$ 衡量当前状态与高层参考之间的偏差，同时考虑动作能量、稳定性、风险和身体约束；$g_B$ 则表示硬约束或软约束。

一个更现实的多目标局部代价可以写成
\[
\mathcal L_B
=
\lambda_r d(x_t,r_t)
+
\lambda_u\|u_t\|^2
+
\lambda_s C_{\mathrm{stability}}
+
\lambda_e C_{\mathrm{energy}}
+
\lambda_risk C_{\mathrm{risk}}
+
\lambda_c C_{\mathrm{constraint}}.
\]
这揭示了 BPCC 与简单 ``执行命令'' 的根本不同：同一个高层目标，在不同身体状态下可能产生完全不同的局部动作。认知内核只表达语义目标，而 BPCC 负责把语义目标转化为当前身体真正能执行的状态轨迹。

这种结构还允许我们把熟练能力解释为局部闭环逐渐稳定化。初次执行某动作时，高层可能给出较密集的控制参考：
\[
r_t^{K}
\]
不断变化；随着 BPCC 掌握稳定的局部动力关系，高层参考可以逐渐稀疏：
\[
\frac{\partial r_t^{K}}{\partial t}
\downarrow.
\]
最后，高层只需要设置一个较宽的目标区间，而 BPCC 自行完成其中的细节。这就是 ``能力向下沉降'' 在控制意义上的表现。

Fast Control 同样要求对实时性有严格认识。高层语言模型通常以离散 token、动作块或语义事件工作，而 BPCC 必须面向连续或高频离散动力过程。在物理机器人中可能是毫秒到几十毫秒，在数字身体中也可能是 GUI 动画、网络回调、系统事件和异步进程。因此 BPCC 不能以 ``每次控制都等待一次完整高层推理'' 的方式实现。它必须具有独立的控制周期：
\[
t_n=n\Delta t_B,
\]
满足
\[
\Delta t_B\ll\tau_h.
\]

外部已有的 PID、LQR、MPC、adaptive control、trajectory optimization 都可以成为 BPCC 局部控制器的具体实现工具，但 BPCC 概念并不等价于其中任意一种算法。BPCC 是一个架构位置：它规定某类问题必须在身体快速局部闭环中解决。具体身体可以选择线性控制、非线性控制、强化学习策略、神经控制器或混合方法。

因此，我们不应把 BPCC 的 ``智能'' 理解成它拥有复杂语言推理，而应理解为：它能够在约束、预测和实时反馈下，让身体以远高于认知环频率的速度完成稳定的局部选择。只有把这层能力从高层认知中剥离，Agent 才能真正拥有一个可以长期使用、逐渐熟练并持续扩张的身体。

\section{Residual Estimation：以残差连接身体与认知}

BPCC 的第三个核心能力是残差估计。预测与控制本身并不能保证模型永远正确，因为现实环境不会承诺服从 BPCC 当前学到的局部动力模型。真正重要的是，系统能否持续知道 ``自己的预测究竟错了多少，以及这种错误是否值得向更高层报告''。因此残差并不是附加诊断量，而是多尺度智能系统中自底向上信息传播的主要驱动力。

设 BPCC 在时刻 $t$ 对下一状态的预测为
\[
\hat x_{t+1}
=
f_\phi(x_t,u_t),
\]
而现实观测得到
\[
x_{t+1}^{obs}.
\]
最基本残差为
\[
\varepsilon_t
=
x_{t+1}^{obs}
-
\hat x_{t+1}.
\]
然而，仅有欧氏差值通常不足以描述复杂身体。我们更需要多通道残差：
\[
\boxed{
\varepsilon_B
=
(
\varepsilon_{\mathrm{ctrl}},
\varepsilon_{\mathrm{SGC}},
\varepsilon_{\mathrm{FCS}},
\varepsilon_{\mathrm{dyn}},
\varepsilon_{\mathrm{safety}}
)
}
\]
其中 $\varepsilon_{\mathrm{ctrl}}$ 表示动作结果与目标的偏差，$\varepsilon_{\mathrm{SGC}}$ 表示世界语义几何预测与实际观察之间的差异，$\varepsilon_{\mathrm{FCS}}$ 表示预测身体影响与实际身体影响之间的差异，$\varepsilon_{\mathrm{dyn}}$ 表示局部动力模型失配，而 $\varepsilon_{\mathrm{safety}}$ 表示与安全包络的偏差。

残差还必须区分瞬时扰动与结构性失配。若一个误差只持续一个控制周期，很可能只是噪声：
\[
\varepsilon_t\neq0,
\qquad
E[\varepsilon_{t:t+T}]\approx0.
\]
但若存在
\[
\left\|
\frac{1}{T}
\sum_{k=t}^{t+T}
\varepsilon_k
\right\|
>\theta_\varepsilon,
\]
则说明模型具有系统性偏差，或者环境进入了新的动力区间。此时应触发不同等级的响应：小残差由 BPCC 局部修正；中等持续残差触发模型适应；高风险或长期残差则转换成高层事件。

因此，可以定义一个残差升级函数：
\[
R_{\mathrm{up}}
=
F(
\|\varepsilon\|,
\dot{\varepsilon},
T_{\mathrm{persist}},
Risk,
Novelty,
GoalRelevance
).
\]
当
\[
R_{\mathrm{up}}<\theta_1
\]
时，
\[
\text{LocalCorrection};
\]
当
\[
\theta_1\leq R_{\mathrm{up}}<\theta_2
\]
时，
\[
\text{LocalAdaptation};
\]
当
\[
R_{\mathrm{up}}\geq\theta_2
\]
时，
\[
\text{SemanticEscalation}.
\]

这正是本书多尺度理论中 ``异常向上'' 的一个基础实现。成熟系统不能把所有数据向上传输，因为那会重新制造工作记忆瓶颈；但也不能把所有异常都留在低层，否则高层永远不知道现实模型已经失效。残差因此承担一种信息选择作用：\textbf{只有现实对现有低层结构的反证，才有资格要求更昂贵的高层认知介入。}

在这个意义上，Residual Estimation 是 BPCC 与整个 AgentOS 之间最重要的因果接口之一。高层真正需要知道的往往不是 ``身体现在每个关节多少度''，而是 ``当前低层模型是否还能可靠完成既定意图''。因此，BPCC 的成熟度可以部分通过残差压缩比来衡量。设原始身体数据率为 $D_B$，真正需要上浮的高层异常信息率为 $D_E$，则：
\[
\eta_{\mathrm{residual}}
=
\frac{D_E}{D_B}.
\]
在正常熟练运行下，理想状态是
\[
\eta_{\mathrm{residual}}\ll1,
\]
但一旦环境发生结构变化，该比例应及时上升。也就是说，BPCC 应同时做到低日常带宽和高异常敏感度。

\section{Adaptive Learning：身体动力模型的在线适应}

一个静态 BPCC 在初始训练分布中可能表现良好，却很难在真实生命周期中长期有效。身体会磨损，设备参数会漂移，网络延迟会改变，工具会被替换，机器人抓取的对象会变化，数字界面的布局和响应速度也会更新。因此，BPCC 必须具有在线适应能力。然而，这种适应必须被严格限定在身体快速动力域中，而不能演变成无边界的主体自修改。

设 BPCC 的局部动力模型参数为
\[
\phi_t,
\]
预测残差为
\[
\varepsilon_t.
\]
最基本的适应形式可以写成
\[
\phi_{t+1}
=
\phi_t
-
\eta_B
\nabla_\phi
\mathcal L_{\mathrm{pred}}
(\varepsilon_t),
\]
其中 $\eta_B$ 为身体层适应速率。但现实系统不会仅依赖即时梯度，因为瞬时残差可能来自噪声。更合理的是引入时间加权残差：
\[
\bar\varepsilon_t
=
(1-\alpha)\bar\varepsilon_{t-1}
+
\alpha\varepsilon_t,
\]
再根据持续性和置信度决定是否更新：
\[
\Delta\phi_t
=
-\eta_B
\,c_t
\,g(\bar\varepsilon_t),
\]
其中 $c_t\in[0,1]$ 表示观测可信度。

这一机制体现了本书此前反复强调的 ``快变量可塑、慢结构受保护'' 原则。BPCC 的内部动力参数可以相对快速适应，而高层长期结构 $\Theta,\Psi,\Omega$ 的变化频率必须远慢于 BPCC：
\[
\tau_{\phi_B}
\ll
\tau_\Theta
\ll
\tau_\Psi
\ll
\tau_\Omega.
\]
如果每一次身体误差都能直接修改 Agent 的长期世界观或身份，那么短期噪声将不断污染慢变量；反过来，如果 BPCC 完全不能适应，所有身体变化又都会被迫升级给高层处理。因此合理设计必须建立明显的时间尺度隔离。

Adaptive Learning 还需要区分至少三类变化。第一类是参数漂移，例如电机摩擦系数改变、鼠标响应速度变化，这类问题适合局部在线参数更新。第二类是动力模式切换，例如机器人从空载进入负载状态、车辆从干燥路面进入湿滑路面，此时更适合使用多模型或模式识别：
\[
m_t
\in
\{M_1,M_2,\ldots,M_K\}.
\]
第三类是真正的结构变化，例如更换机械臂、操作系统接口大幅升级。这种变化可能已经超出 BPCC 当前模型族的表示能力，需要向更高层报告，而不是依靠无休止局部更新掩盖结构失配。

因此可以定义局部适应充分性：
\[
A_t
=
\frac{
\Delta L_{\mathrm{residual}}
}{
Cost_{\mathrm{adapt}}
}.
\]
若连续适应能够稳定降低误差，
\[
A_t>0,
\]
BPCC 保持局部学习；若多次更新后残差仍然持续：
\[
\frac{\partial
\|\bar\varepsilon\|
}{
\partial n_{\mathrm{adapt}}
}
\approx0,
\]
则说明问题不再是普通参数适应，而应触发结构性升级。

这与传统 adaptive control、online system identification 和 residual learning 有明显联系。BPCC 可以借用这些成熟理论，但本书给予它们一个更大的认知系统位置：局部适应不是孤立算法，而是智能结构迁移的一部分。一个最初需要高层认知反复纠正的身体过程，如果逐渐被 BPCC 吸收并稳定处理，就发生了：
\[
ExplicitCognition
\rightarrow
LocalAdaptiveClosure.
\]
因此 BPCC 的在线学习其实是 Agent ``越来越会使用自己的身体'' 的最低层动力基础。

\section{Private Latent State：为什么 BPCC 可以拥有不透明内部状态}

在本书对 SGC、FCS 的定义中，我们强调 Agent Kernel 与 Body Runtime 之间需要可检查、可解释的公共语义界面。但如果进一步要求 BPCC 内部的每一个状态也必须完全采用人类可读的符号表示，将产生严重性能代价。快速身体动力通常具有高维、连续、非线性和强耦合特征，对这些过程进行高频控制时，紧凑潜在状态往往比显式语义图更加高效。因此，BPCC 应被允许维护自己的私有潜在状态：
\[
\boxed{
z_t^{BPCC}
\in
\mathbb R^d.
}
\]

这个 $z_t^{BPCC}$ 可以由历史传感状态、动作和残差递归更新：
\[
z_{t+1}
=
f_z
\left(
z_t,
o_t,
u_t,
\varepsilon_t
\right),
\]
并用于短期预测：
\[
\hat x_{t+1}
=
f_p(z_t,u_t),
\]
以及局部控制：
\[
u_t
=
\pi_B(z_t,r_t).
\]
因此 BPCC 内部完全可以使用 RNN state、state-space model、latent dynamics、world model 或其他高效表示。

然而，``允许不透明 latent'' 不等于 ``允许隐藏语义权力''。我们必须严格区分：
\[
PrivateFastState
\]
和
\[
PublicEmbodiedState.
\]
前者服务于快速计算，后者服务于整个 Agent 的跨模块理解与治理。二者关系可以表示为：
\[
z_t^{BPCC}
\xrightarrow{D_B}
(
SGC_t,FCS_t,C_t^{BPCC}
),
\]
其中 $D_B$ 是语义解码或状态投影函数，$C_t^{BPCC}$ 表示置信度、能力边界和异常状态。

这样设计的理由在于，多尺度智能不要求所有尺度共享一种表示语言。高层工作记忆适合语义密集、关系稀疏、可解释的结构；低层控制适合连续、高维、高带宽的动力 latent。要求二者完全同构反而会破坏多尺度分工。因此我们提出：
\[
\boxed{
DifferentInternalRepresentation,
SharedFunctionalContract.
}
\]

更形式化地说，BPCC 的 latent 不需要满足：
\[
z_t^{BPCC}
\cong
h_t,
\]
但其对外行为必须满足一个功能合同：
\[
\mathcal C_B
=
(
InputContract,
OutputContract,
SafetyContract,
ConfidenceContract
).
\]
只要 BPCC 能够稳定地把公共语义状态映射到有效身体控制，并把异常、能力边界与实际结果重新投射到公共接口，高层就没有必要理解其内部每个 latent 维度。

这与深度表示学习中的 ``latent representation''、控制中的 ``minimal sufficient state'' 和系统辨识中的内部状态模型存在直接联系。但在 AgentOS 中，我们进一步强调了权限边界：私有 latent 可以优化性能，却不能成为一个绕过公共语义接口、直接修改长期认知结构的隐藏通道。

因此，BPCC 的设计原则不是 ``所有东西都要可解释''，也不是 ``内部完全黑箱即可''，而是：
\[
\boxed{
InspectableBoundary
+
OptimizedInternalDynamics.
}
\]
这是未来 AgentOS 工程中非常重要的一条折中原则。越靠近现实高频控制，越允许高效私有表示；越接近跨模块协调、长期记忆和主体治理，越必须转换成稳定、可检查的公共语义状态。

\section{Public Semantic State：BPCC 必须向上层公开什么}

既然 BPCC 可以拥有私有潜在状态，就自然出现另一个问题：它究竟必须向 Body Runtime 和 Agent Kernel 公开什么？答案不应是 ``公开所有传感器数据''，也不应是 ``只返回动作成功/失败''。公开状态必须恰好覆盖高层进行理解、预测、治理和异常处理所需的最小语义充分集。

在当前架构中，公共具身状态首先由
\[
\boxed{
PublicEmbodiedState
=
SGC+FCS
}
\]
构成。SGC 描述当前可行动现实中的实体、空间、关系和认识论状态；FCS 描述身体和现实正在如何功能性影响 Agent。BPCC 则应在这两种公共状态之上补充自己的预测和控制元信息：
\[
\mathcal P_B
=
(
Confidence,
Residual,
CapabilityMargin,
ControlStatus,
PredictionStatus
).
\]

其中，Prediction Confidence 表示短期模型的可靠度：
\[
c_p(t)\in[0,1].
\]
Control Confidence 表示当前控制器对完成局部参考的把握：
\[
c_u(t)\in[0,1].
\]
Capability Margin 表示当前身体距离控制边界还有多少余量。例如：
\[
m_B
=
\min_j
\left(
b_j-g_j(x,u)
\right),
\]
当 $m_B\rightarrow0$ 时，说明动作正在逼近约束边界。Residual 则反映模型与现实之间的不一致。

这些变量的意义在于，高层真正需要的不是低层如何计算，而是：
\begin{statementbox}
CanTheBodyStillReliablyDoWhatIAsked?
\end{statementbox}
例如一个机器人举起物体时，Kernel 不需要每毫秒读取每个关节电流；但它应该知道 ``当前抓取稳定、剩余力矩裕量很低、预测置信度正在下降''。这是一种语义压缩：
\[
RawBodyDynamics
\rightarrow
ActionRelevantSemanticState.
\]

我们可以把公共状态的充分性定义为：对于高层决策变量 $G_K$，若两个低层状态 $z_1,z_2$ 被映射为同一个公共状态 $y_B$，并且它们对当前高层允许动作集合没有显著差异，则这种压缩对于该任务是可接受的。形式上，如果
\[
D(z_1)=D(z_2)=y_B
\]
且
\[
\mathcal A_K(z_1)
\approx
\mathcal A_K(z_2),
\]
那么高层没有必要知道 $z_1$ 与 $z_2$ 的全部内部区别。

这实际上是一种任务相关的 sufficient representation 思想。它使 AgentOS 可以在低层保留高带宽状态，同时让高层获得稳定低带宽接口。公共状态越成功，系统越能够扩展到不同 Body 而不要求 Kernel 重写内部世界模型。

但公共状态也不能过度压缩。如果 BPCC 将所有内部异常都隐藏成一个 ``正常'' 标签，高层就失去了判断模型是否正在逼近失效的能力。因此公开接口必须至少保留：现实语义状态、身体影响状态、预测置信度、控制置信度、能力余量与异常残差。它们共同构成 Body Runtime 对认知内核的 ``可治理表面''。

由此我们得到一个重要设计原则：
\[
\boxed{
PrivateStateMayBeComplex,
PublicContractMustBeStable.
}
\]
这也是不同身体能够被同一个 Agent Kernel 使用的前提之一。

\section{Local State Decoder：从高速 latent 回到可理解身体状态}

私有潜在状态和公共语义状态之间必须存在一个明确的转换机制，这就是 Local State Decoder 的理论位置。BPCC 可以在内部以最适合预测与控制的方式组织状态，但一旦这些状态需要参与 Agent 的显式认知、长期学习或身体治理，就必须转换成稳定语义表示。

设 BPCC 内部状态为
\[
z_t,
\]
Decoder 定义为
\[
D_\psi:
\mathcal Z_B
\rightarrow
\mathcal Y_B,
\]
其中
\[
\mathcal Y_B
=
\mathcal Y_{SGC}
\times
\mathcal Y_{FCS}
\times
\mathcal Y_{meta}.
\]
于是：
\[
y_t^B
=
D_\psi(z_t)
=
(
\hat{SGC}_t,
\hat{FCS}_t,
Meta_t
).
\]

这里必须注意 ``预测状态'' 与 ``现实观测状态'' 的认识论区别。Decoder 从 BPCC latent 解码出的
\[
\hat{SGC}
\]
本质上仍然属于 Predicted 或 Derived，而不能自动升级成 Observed。真正的观察应由传感器重新经过感知链路确认。因此公共状态中必须保留 provenance：
\[
e_t
\in
\{
Observed,
Derived,
Predicted,
Counterfactual
\}.
\]
若忽略这一标签，BPCC 就可能把自身预测偷偷写回世界模型，使 Agent 逐渐生活在自己预测出来的现实中。这将破坏本书一直坚持的
\[
RealityIsUltimateConstraint.
\]

Decoder 的另一个功能是把低层连续状态压缩成高层事件。例如，机械臂内部可能经历数百个小幅震荡：
\[
z_{t:t+N},
\]
高层真正需要的可能只是：
\[
\text{``Grip instability increasing.''}
\]
因此 Decoder 不只是静态向量映射，而可以包含时间窗口上的模式识别：
\[
D_\psi
\left(
z_{t-k:t}
\right)
\rightarrow
EventCSO.
\]
这正是从身体连续动力学生成 Agent-CSO 的关键入口之一。

为了保证 Decoder 不因 BPCC 在线学习而发生不可控语义漂移，可以要求其满足一定的语义稳定性约束。设更新前后的 Decoder 分别为 $D_{\psi_t}$ 与 $D_{\psi_{t+1}}$，对于标准锚点状态集合 $\mathcal Z_{anchor}$，要求：
\[
\mathbb E_{z\sim\mathcal Z_{anchor}}
\left[
d_{\mathrm{sem}}
\left(
D_{\psi_t}(z),
D_{\psi_{t+1}}(z)
\right)
\right]
<
\epsilon_D.
\]
即低层内部模型可以持续适应，但公共语义含义不能随意漂移。

这与传统 system identification 中的 state decoder、表示学习中的 latent-to-observation model 以及机器人状态估计都有联系。然而在 AgentOS 中，Decoder 还有一个额外作用：它维护不同时间尺度之间的语义契约。BPCC 可以 ``用自己的语言思考身体''，但当它与主体交流时，必须回到 Agent 可以持续理解的世界语义。

因此 Local State Decoder 既是技术适配器，也是多尺度认知边界。它保障了：
\[
FastLatentDynamics
\not\Rightarrow
SemanticChaos,
\]
使低层高度优化和高层长期稳定可以同时存在。

\section{MPC 与局部控制目标：BPCC 的统一优化形式}

将快速预测、快速控制、残差估计和在线适应合并起来以后，一个自然的工程实现形式是 Model Predictive Control（MPC）及其更一般化的预测优化框架。这里需要强调，BPCC 并不等于 MPC；MPC 是一个非常适合表达 BPCC 核心思想的数学工具，因为它天然把短期动力预测、约束和滚动控制组合在一起。

在时刻 $t$，设 BPCC 获得当前局部状态
\[
x_t,
\]
高层参考轨迹
\[
r_{t:t+H},
\]
以及身体约束集合
\[
\mathcal C_B.
\]
BPCC 可以求解：
\[
\begin{aligned}
U_t^\ast
=
\arg\min_{U}
\quad&
\sum_{k=0}^{H-1}
\Big[
\lambda_r
d(\hat x_{t+k},r_{t+k})
+
\lambda_u
\|u_{t+k}\|^2
\\
&+
\lambda_s
C_{\mathrm{stability}}
(\hat x_{t+k})
+
\lambda_f
C_{\mathrm{FCS}}
(\hat x_{t+k})
+
\lambda_risk
C_{\mathrm{risk}}
(\hat x_{t+k})
\Big]
\\
&+
\lambda_T
V_f(\hat x_{t+H})
\end{aligned}
\]
subject to
\[
\hat x_{t+k+1}
=
f_\phi(
\hat x_{t+k},
u_{t+k}
),
\]
\[
g_j(
\hat x_{t+k},
u_{t+k}
)
\leq0,
\qquad
j=1,\ldots,m.
\]

求得整个控制序列后，系统并不一次性执行全部 $U_t^\ast$，而只执行第一个动作：
\[
u_t=u_{t|t}^\ast.
\]
随后现实再次被观测，得到
\[
x_{t+1}^{obs},
\]
重新计算残差：
\[
\varepsilon_{t+1}
=
x_{t+1}^{obs}
-
\hat x_{t+1},
\]
再在下一控制周期重新优化。这就是 receding-horizon control。

这套结构与本书的 World--Agent--World 基本思想高度一致：
\[
Predict
\rightarrow
Act
\rightarrow
ObserveReality
\rightarrow
Residual
\rightarrow
Replan.
\]
它最大的理论价值在于，没有任何预测会被永久当成现实。每一个短期未来都只是候选轨迹，每一个实际动作之后都必须重新让现实 ``投票''。

BPCC 的局部优化目标还必须与高层目标保持严格层级区分。高层可能给出：
\[
G_K=
\text{``把杯子安全送到桌面另一侧''},
\]
但 BPCC 求解的是：
\[
\min
\left(
TrajectoryError
+
Energy
+
Instability
+
Risk
\right).
\]
它不重新决定杯子是否应该被送过去，而只决定在给定主体意图下当前身体怎样完成动作。这就是高层 ``决定意义'' 与低层 ``决定实现'' 的边界。

为了表达高层并不需要给出单一路径，我们还可以把 $r_t$ 扩展成预测包络：
\[
\mathcal E_H(t)
=
\left(
\mathcal X_{\mathrm{target}},
\mathcal C_{\mathrm{allowed}},
\mathcal C_{\mathrm{forbidden}}
\right).
\]
BPCC 只需满足
\[
x_{t:t+H}
\in
\mathcal E_H(t),
\]
便可以在包络内部自主寻找更高效轨迹。这样，高层约束的是未来可行区域，而不是微管理低层动作。这正是
\[
SlowGovernance
\neq
FastMicromanagement
\]
在身体控制中的具体含义。

MPC 还可以与 BPCC 的在线适应结合。动力模型
\[
f_{\phi_t}
\]
根据残差不断更新，控制器又基于最新模型求解新的局部轨迹，从而形成：
\[
Prediction
\rightarrow
Control
\rightarrow
Reality
\rightarrow
Residual
\rightarrow
ModelUpdate
\rightarrow
Prediction'.
\]
这是 BPCC 最完整的内部闭环。

当然，实际实现并不要求每个身体都显式在线求解非线性 MPC。对于超高速控制，可能使用近似策略网络、显式 MPC、LQR、PID、trajectory library 或硬件控制器；关键并不是算法名字，而是保持上述功能结构不变：
\[
\boxed{
LocalPrediction
+
ConstrainedAction
+
RealityFeedback
+
ResidualAdaptation.
}
\]

因此，本章最终可以把 BPCC 的完整状态更新写成：
\[
\boxed{
\begin{aligned}
z_{t+1}
&=
F_z(
z_t,o_t,u_t
),\\
\hat X_{t:t+H}
&=
F_p(
z_t,r_t
),\\
u_t^\ast
&=
\Pi_B(
\hat X_{t:t+H},
r_t,
\mathcal C_B
),\\
\varepsilon_{t+1}
&=
\mathcal R(
o_{t+1},
\hat o_{t+1}
),\\
\phi_{t+1}
&=
\mathcal A(
\phi_t,
\varepsilon_{t+1}
).
\end{aligned}
}
\]

其外部则只暴露：
\[
\boxed{
SGC,
FCS,
Residual,
Confidence,
CapabilityMargin,
ControlStatus.
}
\]



\section{SGC 感知控制：身体不仅作用于现实，也主动选择如何读取现实}

前面的快速控制主要讨论了 BPCC 如何把高层控制参考转换成连续或高频的身体动作。然而，如果我们只保留这一类 ``Action Control''，就仍然隐含了一个传统感知系统的假设：现实首先被动进入传感器，随后 Body Runtime 把已经得到的感知结果组织成 SGC，BPCC 只负责在这个世界表示上执行动作。对于真正具有主动身体的 Agent，这个假设并不充分。身体不仅能够改变现实，还能够主动改变自己观察现实的方式。因此，在 BPCC 中还必须存在一类与直接动作控制平行、但认识论性质不同的控制过程：

\begin{statementbox}
PerceptualControl
\end{statementbox}

或者在 SGC 接口层更加具体地称为

\begin{statementbox}
SGCReadOnlyControl.
\end{statementbox}

所谓 ``只读'' 并不意味着控制过程中没有任何物理或数字状态发生变化，而是指该操作的主要意图不是修改当前任务世界中被观察对象的语义状态，而是改变 Agent 获取、选择、组织或验证现实信息的方式。我们可以把身体控制空间划分为两个基本子空间：

\[
\boxed{
\mathcal U_B
=
\mathcal U_B^{act}
\oplus
\mathcal U_B^{obs}
}
\]

其中

\[
u_t^{act}
\in
\mathcal U_B^{act}
\]

表示对现实对象、身体姿态或任务环境产生直接作用的控制，而

\[
u_t^{obs}
\in
\mathcal U_B^{obs}
\]

表示改变观测配置的感知控制。前者主要回答 ``我要怎样改变现实''，后者回答 ``我要怎样看清现实''。

这种区分对于物理身体和数字身体都成立。对于物理机器人，转动摄像头、调整焦距、改变传感器方向、选择深度相机感兴趣区域、改变某区域采样频率，都可以属于感知控制。对于数字身体，同一结构则更加明显：移动观察视口、选择屏幕区域、放大某个界面、选择特定 UI 元素进行只读检查、读取某个对象的 accessibility 属性、查询一个菜单当前可提供但尚未执行的功能、获取光标当前位置附近的语义对象，原则上都不应立即被解释成对任务现实的修改。这些操作首先改变的是

\begin{statementbox}
ObservationState
\end{statementbox}

而不是

\begin{statementbox}
TaskWorldState.
\end{statementbox}

因此，我们需要在身体状态中进一步区分

\[
x_t^B
=
\left(
x_t^{body},
x_t^{obs},
x_t^{world}
\right),
\]

其中 $x_t^{body}$ 表示身体自身状态，$x_t^{obs}$ 表示当前观测配置，而 $x_t^{world}$ 表示通过身体获得的当前外部世界估计。感知控制主要作用于

\[
x_t^{obs},
\]

其动力学可以写成

\[
x_{t+1}^{obs}
=
f_{obs}
\left(
x_t^{obs},
u_t^{obs}
\right),
\]

随后新的观测配置产生新的感知结果：

\[
o_{t+1}
=
\mathcal O
\left(
x_{t+1}^{obs},
World_{t+1}
\right),
\]

再由 Body Runtime 重新生成

\[
SGC_{t+1}.
\]

由此形成一条与动作闭环不同的局部循环：

\begin{statementbox}
\centering
\adjustbox{max width=.96\linewidth}{\(\displaystyle SGC_t
\rightarrow
PerceptualQuestion
\rightarrow
u_t^{obs}
\rightarrow
NewObservation
\rightarrow
SGC_{t+1}.\)}
\end{statementbox}

这条闭环极其重要，因为它使 SGC 从 ``被动生成的场景表示'' 转变成 ``可以被身体主动探索的现实语义空间''。

对于数字身体，我们可以用一个非常简单的例子理解这一点。假设当前 SGC 中存在一个屏幕窗口对象

\[
O_{window},
\]

其中包含多个尚未充分解析的区域。高层 Agent 并不一定需要立即点击任何按钮，它可能只是产生一个认识论请求：

\[
q_t=
\text{``这个区域里面当前有哪些可用操作？''}
\]

BPCC 可以把这个请求转换成只读感知控制：

\[
u_t^{obs}
=
\text{SelectROI}(R),
\]

随后 Body Runtime 对区域 $R$ 进行更高精度读取，并返回

\[
SGC_{t+1}
=
SGC_t
\cup
\{
O_1,O_2,\ldots,O_n
\}.
\]

在这个过程中，世界的任务语义状态可能完全没有改变：

\[
\Delta World^{task}
\approx0,
\]

但 Agent 的可观测结构明显增加：

\[
\Delta I(SGC;World)>0.
\]

因此，感知控制可以被理解为以最小现实扰动换取最大有效信息增益的问题。一个理想的局部感知策略可以写成

\[
u_{t}^{obs*}
=
\arg\max_{u^{obs}}
\left[
\mathbb E
\left(
\Delta I_{task}
\mid
u^{obs}
\right)
-
\lambda_o C_{obs}(u^{obs})
-
\lambda_r R_{sideeffect}(u^{obs})
\right],
\]

其中 $\Delta I_{task}$ 表示该观测动作预计带来的任务相关信息增益，$C_{obs}$ 表示观测成本，而 $R_{sideeffect}$ 表示该所谓 ``只读'' 操作实际上对现实产生副作用的风险。后一项尤其重要，因为现实中的很多感知动作并不存在绝对的只读性。例如摄像头转动改变了机器人姿态；鼠标 Hover 可能触发菜单展开；选择文本可能改变网页的 selection state；读取某些远程资源甚至可能触发服务器日志或访问计数。因此 ``Read-Only'' 在 AgentOS 中应定义成：

\begin{definition*}
不以改变目标世界语义状态为主要目的，并受到副作用约束的观测型动作。
\end{definition*}

而不是声称它在物理意义上绝对不改变任何状态。

从 SGC 的角度看，这类操作至少包括四种基本形式。第一种是 \emph{Spatial Selection}，即对已有 SGC 中的某一区域、对象或关系进行选择，提高其感知分辨率。第二种是 \emph{Semantic Inspection}，即要求 Body Runtime 返回对象已经存在但尚未完全暴露的固定属性、可供性或功能描述。第三种是 \emph{Viewpoint Control}，即通过视口、相机、光标位置或数字身体位置改变观察原点。第四种是 \emph{Epistemic Verification}，即对一个 Predicted 或 Derived SGC 结构执行现实查询，把它重新验证为 Observed，或者发现预测错误。

因此我们可以定义

\[
u_t^{obs}
=
\left(
u_t^{roi},
u_t^{inspect},
u_t^{view},
u_t^{verify}
\right).
\]

它们共同构成 BPCC 的 SGC 感知控制子空间。

这一设计还使 ``数字身体的位置'' 获得了更加严格的工程含义。对于桌面 Agent，鼠标、焦点、当前窗口、viewport、selection 等并不只是 UI API 的零散状态，而构成数字身体在 SGC 中的局部姿态：

\[
Pose_{digital}
=
(
Cursor,
Focus,
Viewport,
Selection,
ActiveSurface
).
\]

因此，选择一个区域并不是抽象软件命令，而类似于物理机器人 ``把眼睛和手的注意中心移动到某个位置''。SGC 应围绕这一数字身体位置重新组织距离、可达性和操作关系。例如某个按钮距离当前 cursor 很近，而另一个对象位于另一个窗口；这些关系都可以成为局部控制策略的一部分。

这里尤其需要区分三个容易混淆的概念：

\[
\boxed{
Attention
\neq
PerceptualControl
\neq
WorldAction.
}
\]

Attention 是 Agent 内部决定哪些已有结构值得获得更多认知资源；Perceptual Control 是通过身体主动改变下一步获得什么现实信息；World Action 则是试图改变现实状态本身。三者可以连续连接：

\[
AttentionRequest
\rightarrow
PerceptualControl
\rightarrow
NewSGC
\rightarrow
Decision
\rightarrow
WorldAction,
\]

但不能在理论上合并成同一操作。

这一区分对工作记忆有限容量理论也具有直接意义。高层不应该一次性获得整个高分辨率世界，而可以首先保持一个粗粒度 SGC，然后仅在必要时向 BPCC 发出感知查询：

\[
SGC^{coarse}
\xrightarrow{Query}
SGC^{local\ fine}.
\]

因此：

\[
\boxed{
PerceptualControl
=
ActiveAdmissionControlBeforeWorkingMemory.
}
\]

也就是说，CCM 与 Boundary 主要管理 ``什么结构进入认知''，而 BPCC 的感知控制更早一步决定 ``身体下一刻去获取什么结构''。二者共同构成从现实到工作记忆的主动信息治理链。

这一点也解释了为什么真实的 SGC 不应该只是一次性的视觉编码。SGC 必须支持可查询性：

\[
Query(SGC_t,q)
\rightarrow
PerceptualAction
\rightarrow
Observation
\rightarrow
SGC_{t+\Delta}.
\]

在数字环境中，Agent 可以选择一个菜单但不执行其中任何命令，只要求系统返回该菜单包含哪些现实功能；也可以询问某个 UI 对象目前是否 enabled、其操作会影响哪个对象、当前窗口是否具有写权限。这类信息属于现实对象的可供性结构，应重新进入 SGC，而不是由高层语言模型凭经验猜测。

最后，感知控制和直接控制必须在 SGC 中形成不同的认识论记录。可以为每个 Body Action 加入

\[
ActionMode
\in
\{
OBSERVE,
ACT,
SAFETY
\}.
\]

其中

\[
OBSERVE
\]

表示认识论型身体操作；

\[
ACT
\]

表示现实修改型控制；

\[
SAFETY
\]

表示紧急保护型控制。

如果执行

\[
u_t^{obs},
\]

SGC 更新主要体现为

\[
\Delta SGC^{knowledge},
\]

而如果执行

\[
u_t^{act},
\]

则必须等待现实重新观测以后确认

\[
\Delta SGC^{world}.
\]

特别需要禁止：

\[
u_t^{act}
\Rightarrow
SGC_{t+1}^{Observed}
\]

这种直接写入，因为 ``已经发出动作'' 并不等价于 ``现实已经发生变化''。正确链路必须是：

\[
u_t^{act}
\rightarrow
World
\rightarrow
Sensor
\rightarrow
ObservedSGC_{t+1}.
\]

相反，控制器可以暂时生成：

\[
SGC_{t+1}^{Predicted},
\]

但必须带有明确的 epistemic provenance，等待真实观测校验。

因此 BPCC 的完整控制概念不应该只写成 ``身体动作控制器''，而应该写成：

\[
\boxed{
BPCCControl
=
PerceptualControl
+
WorldActionControl
+
LocalSafetyResponse.
}
\]

三者分别解决：

\begin{statementbox}
我下一步应该看哪里？
\end{statementbox}

\begin{statementbox}
我下一步应该怎样作用于世界？
\end{statementbox}

以及：

\begin{statementbox}
发生危险时身体必须立即做什么？
\end{statementbox}

这三种身体能力与 Fast Prediction、Residual Estimation 和 Adaptive Learning 共同构成一个真正完整的 Body Predictive Control Core。


从本书整个智能动力学体系来看，BPCC 的真正意义并不是为 Agent 加一个运动控制模块，而是建立第一条真正清晰的多尺度认知边界：\textbf{高层负责意义、目标与长期结构，低层负责高速预测、局部闭环和身体适应；高层不微观管理低层，低层也不夺取主体权力，而二者通过现实反馈保持持续耦合。}

当这一结构成立以后，一个 Agent 才真正从 ``拥有一个动作接口'' 跨越到 ``拥有一个能够被长期学习、逐渐熟练并持续适应的身体''。这也是 BPCC 在完整 AgentOS 架构中的理论位置。
